\subsection{Fuerza Bruta / Simulación Directa (paradoja del cumpleaños)}
\label{alg:7-2}

\graybold{Preguntas guía:}

\begin{itemize}
\item ¿El tamaño de la entrada es pequeño (N ≤ \ten{3} o similar) y no hay una fórmula cerrada evidente?
\item ¿Puedo calcular la respuesta iterando directamente sobre el espacio del problema sin que se dispare el tiempo?
\item ¿El problema es más de matemáticas/probabilidad que de estructuras de datos?
\end{itemize}

\graybold{Utilidad:} Reconocer cuándo NO hace falta un algoritmo sofisticado: con restricciones pequeñas, iterar directamente (simulando la definición del problema) es la solución más simple, rápida de programar y menos propensa a errores. Identificar esto ahorra tiempo valioso en el contrarreloj.

\graybold{Complejidad:} Depende del problema; en este caso \bigO{n} con aritmética de punto flotante.

\graybold{Fundamento teórico:} Se calcula directamente la probabilidad de que NO haya coincidencia de cumpleaños entre k personas, multiplicando factores (1 − (k−1)/n) sucesivamente (extensión del problema del cumpleaños), deteniéndose apenas esa probabilidad cae por debajo de 0.5 (la probabilidad de coincidencia supera el 50\%).

\graybold{Ejercicio aplicado}

Dado N (número de días del año), hallar el mínimo número de personas necesarias en una sala para que la probabilidad de que dos compartan cumpleaños supere el 50\%.

\graybold{I/O:} Entrada de un solo entero → readline (patrón 1), no requiere nada especial.

\graybold{Código solución}

\begin{lstlisting}[language=Python]
import sys
input = sys.stdin.buffer.readline
 
n = int(input())
prob_no_coincidencia = 1.0
for k in range(1, n + 2):
    prob_no_coincidencia *= (n - (k - 1)) / n
    if prob_no_coincidencia < 0.5:
        print(k)
        break
\end{lstlisting}
